% Generated by manuscript/build_manuscript.py.
\section{Conclusion}
\label{sec:conclusion}


Varied-horizon forecasting requires one model to serve different request endpoints without changing predictions for their shared future steps. We formulated this requirement as cross-horizon prefix consistency (CHPC) and identified future-region sharing-demand heterogeneity as a complementary decoder-side challenge. To address both, ISCF constructs a scope-indexed forecast field through region-wise generation at multiple output-side sharing extents and integrates its Scope-conditioned Forecasts through Target-Adaptive Allocation. BSCA supports the joint learning of the fused trajectory and all scope lines without changing the inference graph; the resulting horizon-agnostic trajectory provides the nested prefixes returned for different horizon requests.

Across seven multivariate benchmarks and four forecast horizons, ISCF-BSCA used one unified model per dataset, achieved stronger aggregate forecasting accuracy than the evaluated horizon-specific baselines and ranked first under the one-model-all-horizons comparison. The unified architecture also achieved a favorable balance between forecasting accuracy, checkpoint storage and inference memory. Controlled ablations supported the contribution of multi-scope generation, scope-specific projections, Target-Adaptive Allocation and BSCA. A selected internal diagnostic illustrated distinct scope signals and region-dependent allocation, while experiments with DLinear and PatchTST supported decoder compatibility across the two evaluated Encoder families. Together, these results support output-side sharing granularity as a practical design axis for unified forecasters that serve multiple horizons through one prefix-consistent trajectory.
